\documentclass[11pt]{article}

% This is successor documentation.  It is intentionally not included in the
% reconstructed 1995 manual, whose historical text and pagination are frozen.
\usepackage[letterpaper,margin=0.82in]{geometry}
\usepackage[T1]{fontenc}
\usepackage{lmodern}
\usepackage{microtype}
\usepackage{amsmath,amssymb}
\usepackage{array,booktabs,longtable,tabularx}
\usepackage{enumitem}
\usepackage{xcolor}
\usepackage{listings}
\usepackage{fancyhdr}
\usepackage[hidelinks]{hyperref}

\definecolor{navy}{HTML}{17365D}
\definecolor{bluegray}{HTML}{475569}
\definecolor{paleblue}{HTML}{EEF5FB}
\definecolor{paleyellow}{HTML}{FFF8E7}
\definecolor{palered}{HTML}{FFF0F0}
\definecolor{codegray}{HTML}{F5F7FA}

\hypersetup{
  pdftitle={TAORYX Extensions and Verification Guide},
  pdfauthor={TAORYX project},
  pdfsubject={Fidelity modes, controller synthesis, trajectory scoring, and verification}
}

\setlist{nosep,leftmargin=*}
\setlist[itemize]{label=\textbullet}
\setlength{\parindent}{0pt}
\setlength{\parskip}{5pt}
\renewcommand{\arraystretch}{1.18}
\newcolumntype{Y}{>{\raggedright\arraybackslash}X}
\newcommand{\taoryx}{TAORYX}
\newcommand{\pathcode}[1]{\path{#1}}
\newcommand{\warningbox}[1]{%
  \par\noindent\colorbox{paleyellow}{\parbox{0.94\linewidth}{#1}}\par
}
\newcommand{\claimbox}[1]{%
  \par\noindent\colorbox{paleblue}{\parbox{0.94\linewidth}{#1}}\par
}

\lstdefinestyle{taoryxcode}{
  basicstyle=\ttfamily\small,
  backgroundcolor=\color{codegray},
  frame=single,
  rulecolor=\color{bluegray!35},
  columns=fullflexible,
  keepspaces=true,
  breaklines=true,
  showstringspaces=false,
  tabsize=2
}
\lstset{style=taoryxcode}

\pagestyle{fancy}
\fancyhf{}
\fancyhead[L]{\small\textsc{TAORYX extensions and verification}}
\fancyhead[R]{\small\thepage}
\renewcommand{\headrulewidth}{0.4pt}

\title{\textcolor{navy}{TAORYX Extensions and Verification Guide}\\
  \large Fidelity modes, plant data, controller synthesis, and trajectory evidence}
\author{TAORYX project}
\date{Successor documentation --- July 2026}

\begin{document}

\maketitle
\thispagestyle{empty}

\claimbox{\textbf{Scope and claim boundary.} This guide documents successor-side
\taoryx{} behavior. It is not a revision of the reconstructed TAOS User's Manual
and does not establish TAOS 96.0 behavioral compatibility. The historical manual
remains the authority for documented 1995 syntax; these extensions are explicitly
labelled \texttt{taoryx-extension} unless a separate source claim is recorded.}

This PDF is the narrative design and verification guide. The canonical
extension Markdown pages are the operational reference for exact problem-file
syntax, command examples, and extension-specific contracts; the composite
build packages both views together. When a detail changes, update the
operational page and its tests first, then keep this guide's rationale and
claim boundary aligned with it.

The manual-parallel language reference is maintained separately in
\pathcode{docs/latex/taoryx_language_reference.tex}. It owns the successor
mathematics, problem grammar, and table-file syntax sections that mirror the
historical manual's Methods, Problem Files, and Table Files organization.

\tableofcontents
\newpage

\section{Purpose and operating principles}

The practical goal is to make a new vehicle and a new controlled trajectory
repeatable without hiding physical assumptions in a bespoke Python runner. A
vehicle enters through catalog metadata, verified source tables, a plant adapter,
and ordinary problem files. The same harness then exercises the vehicle from a
low-fidelity model to a full rigid-body model.

The workflow has four governing principles:

\begin{enumerate}
  \item \textbf{One canonical plant contract.} A vehicle adapter owns frames,
    units, tables, mass properties, propulsion, actuators, and force/moment
    evaluation. Integrators consume the resulting canonical SI state and wrench.
  \item \textbf{Progressive fidelity.} Begin with true 3-DOF translation, add a
    kinematic attitude bridge, and only then release six-degree-of-freedom
    attitude and rate dynamics.
  \item \textbf{Plant truth before controller tuning.} Trimming, local
    linearization, control-direction probes, and table-margin checks precede any
    waypoint or mission score.
  \item \textbf{Gates and scores are different.} A required validity gate can
    block a claim. A quality score describes how much margin a passing objective
    retained; it cannot turn a failed gate green.
\end{enumerate}

The main authoring locations are:

\begin{itemize}
  \item \pathcode{verification/vehicle_models.yaml}: SI vehicle and source
    contract;
  \item \pathcode{verification/vehicle_catalog.yaml}: public vehicle and table
    discovery for the common harness;
  \item \pathcode{verification/trim_specs.yaml}: bounded operating points and
    plant residual names;
  \item \pathcode{verification/controller_designs.yaml}: controller channels,
    allocation, and design method;
  \item \pathcode{verification/lqr_scaling_profiles.yaml}: gentle, standard,
    and aggressive controller profiles;
  \item \pathcode{verification/fidelity_ladder.yaml}: matched 3-DOF, bridge,
    and 6-DOF scenario products;
  \item \pathcode{examples/}: generated or hand-authored ordinary \texttt{.prb}
    examples, never hidden runner inputs.
\end{itemize}

\section{Successor language reference}
\label{sec:language-reference}

The focused operational reference for the successor language is the Markdown
page \pathcode{docs/extensions/} \pathcode{taoryx-language-reference.md}.
This section keeps
the same surface visible in the buildable LaTeX guide: exact syntax belongs in
the Markdown reference and its fixtures, while this guide records the AST and
runtime claim boundary.

Select the successor profile explicitly when parsing an extension file:

\begin{lstlisting}
from taoryx.language import GrammarProfile, parse_problem_file

document = parse_problem_file("mission.prb", profile=GrammarProfile.TAORYX)
\end{lstlisting}

The default \texttt{taos96} profile remains unchanged. It rejects successor
directives rather than silently interpreting them as historical syntax.

\begin{longtable}{p{0.25\linewidth}p{0.17\linewidth}p{0.47\linewidth}}
\toprule
\textbf{Surface} & \textbf{Scope} & \textbf{Typed and runtime meaning}\\
\midrule
\endhead
\pathcode{*3dof}, \pathcode{*6dof}, \pathcode{*sixdof} & problem &
  \texttt{DofDirectiveBlock}; aliases normalize to point-mass or rigid-body 6-DOF.\\
\pathcode{*mode ...} & problem &
  \texttt{ModeBlock}; explicit dynamics mode.\\
\pathcode{*method ...} & problem &
  \texttt{ExtensionBlock}; method and options are preserved.\\
\pathcode{*mass ...}, \pathcode{*ptmass} & segment &
  \texttt{ExtensionBlock}; mass-model transition.\\
\pathcode{*deployed ...} & trajectory &
  \texttt{ExtensionBlock}; deployment initialization.\\
\pathcode{*cases} & problem &
  \texttt{ExtensionBlock}; case columns and rows for batch adapters.\\
\pathcode{surface\_ref}, \pathcode{surface\_azm}, \pathcode{surface\_dist} & define / expr. &
  Typed helper statements and expression AST calls.\\
\pathcode{*runtime ...} & problem &
  \texttt{RuntimeBlock}; shared runtime contract declarations.\\
\pathcode{aero\_force}, \pathcode{aero\_moment}, \pathcode{inertia} & tables &
  Typed successor table documents.\\
\bottomrule
\end{longtable}

The aliases are source-preserving: \pathcode{*3dof} maps to
\texttt{point-mass}, while \pathcode{*6dof} and \pathcode{*sixdof} map to
\texttt{rigid-body-6dof}. Parser output includes the lossless source, typed
semantic AST, and source-located diagnostics. A successful parse is therefore
not equivalent to a runtime-complete case; adapters may return an explicit
unsupported-feature or unavailable-input result.

The smallest focused corpus is under
\pathcode{examples/taoryx/syntax\_fragments/}; combined examples are under
\pathcode{examples/taoryx/full\_examples/}. The same corpus can be regenerated
with parser reports, normalized artifacts, telemetry, and plots:

\begin{lstlisting}[language=bash]
PYTHONPATH=src .venv/bin/python examples/run_corpus.py \
  --family taoryx --execute --output artifacts/examples/taoryx
\end{lstlisting}

The local runtime executes the explicitly supported successor subset. Neither
this guide nor the Markdown reference claims bit-for-bit behavior of the
historical TAOS 96 executable or its complete table library.

\section{The fidelity ladder}
\label{sec:fidelity}

The scenario contract names three tiers. The contract uses the stable names
\texttt{3dof}, \texttt{pseudo\_6dof}, and \texttt{6dof}; problem files use the
current mode directives shown below.

\begin{longtable}{p{0.19\linewidth}p{0.27\linewidth}p{0.45\linewidth}}
\toprule
\textbf{Tier} & \textbf{Problem-file mode} & \textbf{Physical content and evidence meaning}\\
\midrule
\endhead
True 3-DOF & \pathcode{*mode point-mass} &
  Translational position and velocity are integrated. Aerodynamic angles,
  bank, thrust direction, and applicable controls are prescribed or solved by
  the point-mass guidance layer. This is a vehicle reduction only when it uses
  the same atmosphere, tables, propulsion, mass, and command history as the
  higher-fidelity cases.\\
Kinematic 3+3 bridge & \pathcode{*mode kinematic-6dof} &
  Translation remains the point-mass plant while attitude is prescribed,
  scheduled, or lagged. This is useful for frame, attitude-command, force
  direction, and controller development. It is not free rotational dynamics;
  its overlapping position and speed curves are not independent evidence.\\
Full rigid-body 6-DOF & \pathcode{*mode rigid-body-6dof} &
  Quaternion attitude, body rates, inertia, moments, actuators, allocation,
  propulsion moments, and applicable mass/event dynamics are integrated. This
  is the tier required for a full rigid-body plant or controller claim.\\
\bottomrule
\end{longtable}

Minimal examples are ordinary problem files:

\begin{lstlisting}
*mode point-mass
*mode kinematic-6dof
*mode rigid-body-6dof
\end{lstlisting}

The bridge is intentionally a development rung. It can expose an incorrect
body-to-world rotation or a wrong alpha/beta sign before the full attitude
solver is introduced. It must report prescribed versus achieved attitude,
attitude lag, body rates, and force-direction consistency; plotting the same
altitude three times is not a fidelity demonstration.

\subsection{Matched parity versus fidelity separation}

There are two valid experiments, and they must not be mixed.

\begin{description}[leftmargin=1.45in,labelwidth=1.25in,style=sameline]
  \item[Reduction parity] Use identical initial state, mass, environment,
    tables, propulsion, commands, events, duration, integrator, and termination
    policy. Constrain the 6-DOF attitude and rates to the 3-DOF assumptions.
    Compare initial $\dot V$, $\dot h$, $\dot\gamma$, and $\dot\chi$ before
    comparing histories.
  \item[Fidelity separation] Release attitude, rates, actuator lag, allocation,
    and moments. The resulting six-axis trajectory is not expected to overlay
    the 3-DOF curve. Its success criteria are plant/source agreement, closure,
    convergence, bounded controls, and family-appropriate mission behavior.
\end{description}

Every run carries a \texttt{ScenarioContract}. The overlay tool must refuse a
comparison when any physical input hash differs; only the dynamics tier may
differ in a reduction-parity overlay. This prevents a unit conversion or a
different vehicle mass from being mistaken for a dynamics effect.

\section{Vehicle data required by each mode}
\label{sec:data-contract}

The amount of data grows with the physical effects being released. A missing
source value is not silently replaced by a convenient constant; it is recorded
as a warning, diagnostic, or block in the onboarding and plant reports.

\begin{longtable}{p{0.21\linewidth}p{0.23\linewidth}p{0.45\linewidth}}
\toprule
\textbf{Data} & \textbf{3-DOF / bridge} & \textbf{Full 6-DOF addition}\\
\midrule
\endhead
Canonical state & Position, velocity, epoch, and local navigation frame. &
  Quaternion body-to-world convention, body rates, CG location, and a declared
  state ordering.\\
Units and geometry & Mass, reference area, and any lift/drag reference. &
  $S$, $b$, $c$, CG and aerodynamic moment reference, thrust stations, and
  SI-converted inertia tensor.\\
Environment & Gravity, atmosphere, wind, and geodetic/ECEF conversion. &
  The same environment, plus gravity projection into the body frame and all
  rotating-frame derivative conventions used by the rigid-body RHS.\\
Aerodynamics & $C_L$, $C_D$, side force, angle ranges, bank/force direction. &
  Six body or wind-axis force coefficients, $C_l$, $C_m$, $C_n$, rate
  derivatives, control increments, query margins, and out-of-envelope policy.\\
Propulsion & Thrust magnitude and mass flow when active. &
  Body thrust vector, thrust moment, actuator or throttle dynamics, fuel flow,
  and event behavior.\\
Mass model & Mass and, if changing, a continuous mass-flow rule. &
  $m(t)$, CG migration, inertia update, propellant tanks, separation/jettison
  events, and parallel-axis bookkeeping.\\
Controls & Prescribed angle, bank, thrust, or reduced control command. &
  Command/achieved positions, rates, limits, lag, saturation, allocator, and
  signed direction probes for every control channel.\\
\bottomrule
\end{longtable}

For a six-axis source model, the minimum dimensionalization contract is:

\begin{align}
  \boldsymbol F_{A,B} &= q S\begin{bmatrix}C_X&C_Y&C_Z\end{bmatrix}^{T},\\
  \boldsymbol M_{A,B} &= q S\begin{bmatrix}b C_l&c C_m&b C_n\end{bmatrix}^{T},\\
  q &= \tfrac{1}{2}\rho V_{\mathrm{rel}}^2.
\end{align}

The reference dimensions, coefficient normalization, axis ordering, and table
clamp/extrapolation policy belong to the source record. The harness checks these
before it runs a trajectory.

\subsection{How to read the vehicle data contract}

The data table is not a shopping list of optional tuning constants. Each entry
answers a physical question that becomes active when a fidelity tier releases
that effect:

\begin{description}[leftmargin=1.55in,labelwidth=1.35in,style=sameline]
  \item[State and frames] What does each state component mean, in which frame
    is it stored, and which direction does a positive rotation or force use?
  \item[Reference geometry] Which area, span, chord, CG, and moment reference
    were used to nondimensionalize the source coefficients?
  \item[Aerodynamic model] Which coefficient names, axes, interpolation rules,
    and valid ranges produce the force and moment queried by the plant?
  \item[Mass properties] What are $m$, CG, and $I$ at this operating point, and
    what explicit provider updates them as fuel, payload, or stages change?
  \item[Effectors] Which command produces which signed force or moment, through
    which actuator, lag, limit, and allocation rule?
  \item[Provenance] Which source file, page, hash, digitization method, and
    uncertainty basis support each estimate?
\end{description}

For example, the following metadata declares a small fixed-wing reference
point, but does not by itself claim that the aerodynamic coefficients are
truth:

\begin{lstlisting}
*runtime status vehicle
  reference-area=16.0 reference-span=10.0 reference-length=1.6
  mass-kg=1200 dry-mass-kg=950
  inertia-x=4200 inertia-y=5100 inertia-z=8600
  inertia-mass-model=linear-dry-mass
  inertia-x-dry=3900 inertia-y-dry=4700 inertia-z-dry=7800
  aero-source-quality=estimated aero-uncertainty-fraction=0.15
*runtime actuator elevator minimum=-0.35 maximum=0.35
*runtime actuator rudder minimum=-0.30 maximum=0.30
\end{lstlisting}

The corresponding source record must still identify the coefficient table axes
and their units. If the table uses $\alpha$ in degrees but the state uses
radians, the adapter converts before the lookup and records that conversion.
If the source reports $C_D$ and $C_L$ in wind axes, the adapter transforms the
wind-axis force into the body frame before combining it with propulsion and
gravity. Never place a coefficient named $C_D$ directly into a body-$x$ force
without checking the source convention.

\subsection{Worked data dimensionalization example}

Suppose a source-backed fixed-wing case declares

\begin{center}
\begin{tabular}{llllll}
\toprule
$m$ & $V$ & $\rho$ & $S$ & $b$ & $c$\\
\midrule
1200 kg & 60 m/s & 1.225 kg/m$^3$ & 16 m$^2$ & 10 m & 1.6 m\\
\bottomrule
\end{tabular}
\end{center}

The dynamic pressure is
\[
  q=\tfrac12(1.225)(60)^2=2205\ {\rm Pa}.
\]
In straight, level flight the first estimate of required lift is $mg$, so

\[
  C_{L,\mathrm{required}}
  \approx \frac{1200(9.81)}{2205(16)}=0.334.
\]

If the source coefficient at the trimmed angle is $C_L=0.334$, the lift
balance is plausible. If $C_D=0.030$, its dimensional drag is approximately
$1058$ N, so the first propulsion estimate is a thrust of that order, before
including climb angle, propulsion inclination, or other force components. These
are pre-solve checks, not a replacement for the full residual evaluator. They
catch a factor-of-1000 density error, a square-degree/radian mistake, or a
wrong reference area before an optimizer hides it in a control command.

The same discipline applies to moments. For a source moment reference point,
first compute
\[
  \boldsymbol M_{\mathrm{ref}}
  =qS\begin{bmatrix}bC_l&cC_m&bC_n\end{bmatrix}^{T}.
\]
If $\boldsymbol r_{\mathrm{ref}\to CG}$ is the vector from the aerodynamic
reference to the CG, transfer the force-induced moment explicitly:
\[
  \boldsymbol M_{CG}
  =\boldsymbol M_{\mathrm{ref}}
   -\boldsymbol r_{\mathrm{ref}\to CG}\times\boldsymbol F_{A}.
\]
The sign follows from the declared vector direction; the adapter should test
the transfer with a simple off-axis force. Propulsion uses the same rule,
\[
  \boldsymbol M_{T,CG}
  =\boldsymbol M_{T,local}
   +\boldsymbol r_{CG\to T}\times\boldsymbol F_T,
\]
where the vector is explicitly from CG to the thrust application point.

\subsection{Data validation before trim}

A new vehicle should not begin with LQR weights. It should produce a small
data packet containing:

\begin{enumerate}
  \item a source-to-SI table with every unit conversion and reference quantity;
  \item a frame diagram or written basis definition for body, wind, world,
    gravity, and thrust vectors;
  \item coefficient names, table axis order, valid bounds, interpolation policy,
    and at least one in-range and boundary query;
  \item mass, CG, inertia, dry-mass, and mass-property-provider behavior at
    every intended operating point;
  \item control direction probes showing the sign of each force and moment
    response, including actuator limits and achieved-command telemetry; and
  \item an uncertainty note distinguishing measured, digitized, estimated, and
    assumed aerodynamic quantities.
\end{enumerate}

The onboarding report should point to each missing item. A visually plausible
trajectory is not a substitute for this packet: a wrong moment reference or
body-axis sign can still produce a stable-looking trajectory with the wrong
physics.

\subsection{What a new vehicle author supplies}

The generic onboarding path is deliberately small:

\begin{enumerate}
  \item preserve the source bundle and record its provenance and model path;
  \item add one SI vehicle entry and select a family contract;
  \item bind every verified \texttt{.tbl} file and add one golden catalog case;
  \item add a metadata problem-generation profile; generated \texttt{.prb}
    files are products, not the source of truth;
  \item add a bounded trim specification and controller design;
  \item add signed control-direction probes;
  \item add gentle, standard, and aggressive LQR profile entries;
  \item run the onboarding report before running a mission.
\end{enumerate}

\begin{lstlisting}
python tools/validate_vehicle_onboarding.py --vehicle new_vehicle
python tools/dev.py generate-problems
python tools/dev.py vehicles
python tools/dev.py control-directions
python tools/auto_tune_controllers.py --vehicle new_vehicle
\end{lstlisting}

The report points to the exact catalog path, diagnostic code, missing file or
field, and repair hint. A vehicle reaches \texttt{ready} before it can be
considered for a plant-golden claim. A source-only trim remains a warning or
block until its residual adapter exists.

\section{Trim, local derivatives, and the controller architecture}
\label{sec:controllers}

The controller stack is separated from the plant adapter:

\begin{enumerate}
  \item \textbf{Plant and convention firewall:} ingest tables, resolve units,
    construct frames, verify quaternion identity, check alpha/beta signs, and
    evaluate force/moment directions.
  \item \textbf{Bounded trim:} solve the source-backed equilibrium with the
    exact evaluator used for propagation.
  \item \textbf{Local linearization:} obtain true state-derivative Jacobians at
    the accepted trim, not an unlabelled force or moment Jacobian.
  \item \textbf{Controller synthesis:} build LQR, rate, attitude, and outer
    guidance layers through an explicit control contract and allocator.
  \item \textbf{Mission evaluation:} run the fixed problem file, preserve raw
    telemetry, and score objectives only after plant gates pass.
\end{enumerate}

\subsection{Plant-bound trim}

Each trim entry names state variables, control variables, residual variables,
bounds, initial guesses, and residual scales. The vehicle-specific adapter is
the function
\[
  r = \operatorname{plant\_residual}(x,u),
\]
which must call the same source tables, frames, actuators, propulsion, and mass
model used during propagation. The common solver uses bounded nonlinear least
squares and retains both scaled and unscaled residuals.

For a fixed-wing equilibrium, a typical residual vector is
\[
  r_{\mathrm{trim}} =
  \begin{bmatrix}
    \dot u & \dot v & \dot w & \dot p & \dot q & \dot r
  \end{bmatrix}^{T}.
\]
For a rotorcraft, the residual may instead include translational acceleration,
angular acceleration, and the rotor/motor equilibrium. The residual names and
units are part of the trim artifact.

The trim result records the solved state, solved controls, residual components,
solver status, iteration count, bounds, and acceptance tolerance. A numerical
solver success flag is not a plant success flag: table margins, saturation, and
independent closure remain separate gates.

\subsection{What a trim means}

A trim is one declared operating point, not a universal vehicle state. The
author first fixes the external conditions and targets: for example, mass,
altitude, atmosphere, wind, speed, flight-path angle, heading, and whether the
case is straight flight, a coordinated turn, a hover, or a powered climb. The
unknowns then contain only the state and controls that the adapter is allowed
to change. Everything else is held by the trim specification or derived from
it.

For a straight, level fixed-wing example, a useful local choice is
\[
  x_{\mathrm{trim}}=[\alpha,\beta,p,q,r]^{T},\qquad
  u_{\mathrm{trim}}=[\delta_e,\delta_a,\delta_r,\delta_T]^{T},
\]
with target $V$, altitude, and $\gamma=0$. The body velocity is reconstructed
from the air-relative speed and angle convention. In one common convention,
\[
  u=V\cos\alpha\cos\beta,\qquad
  v=V\sin\beta,\qquad
  w=V\sin\alpha\cos\beta.
\]
The source convention may differ, especially for projected versus Euler
sideslip, so the adapter must use the repository's declared angle contract.
The trim controls are not guessed from a waypoint. They are solved so that
the forces, moments, and any target kinematic conditions close at this one
operating point.

\subsection{Body-frame trim equations}

Let $\boldsymbol v_B=[u,v,w]^T$ be body velocity, let
$\boldsymbol\omega_B=[p,q,r]^T$ be body angular rate, and let all forces and
moments below be resolved about the CG in body axes. The translational Newton
equation used by the residual adapter is
\[
  m\left(\dot{\boldsymbol v}_B
    +\boldsymbol\omega_B\times\boldsymbol v_B\right)
  =\boldsymbol F_{A,B}+\boldsymbol F_{T,B}+\boldsymbol F_{g,B},
  \qquad
  \boldsymbol F_{g,B}=R_{I\to B}\,m\boldsymbol g_I.
\]
The rotational Euler equation is
\[
  I\dot{\boldsymbol\omega}_B
  +\boldsymbol\omega_B\times(I\boldsymbol\omega_B)
  =\boldsymbol M_{A,B}+\boldsymbol M_{T,B}+\boldsymbol M_{\mathrm{control},B}.
\]
For diagonal principal inertia,
$I=\operatorname{diag}(I_x,I_y,I_z)$; for a source tensor, retain the full
tensor and its declared frame. The residuals should be formed as left side
minus right side:
\begin{align}
  \boldsymbol r_F &\equiv
    m\left(\dot{\boldsymbol v}_B+\boldsymbol\omega_B\times\boldsymbol v_B\right)
    -\boldsymbol F_{\mathrm{total},B},\\
  \boldsymbol r_M &\equiv
    I\dot{\boldsymbol\omega}_B
    +\boldsymbol\omega_B\times(I\boldsymbol\omega_B)
    -\boldsymbol M_{\mathrm{total},B}.
\end{align}

For straight, level, constant-speed flight with zero body rates, these reduce
to $\boldsymbol F_{\mathrm{total},B}=0$ and
$\boldsymbol M_{\mathrm{total},B}=0$. For a coordinated turn or a hover with
nonzero rates, do not force those terms to zero: retain the target acceleration
and angular-rate terms in the residual. This distinction is why a trim adapter
must know the maneuver condition rather than merely set every derivative to
zero.

The force and moment assembly follows a fixed sequence:

\begin{enumerate}
  \item Compute relative air velocity from vehicle velocity and wind, then
    derive the source-defined $\alpha$, $\beta$, Mach, and dynamic pressure.
  \item Query each aerodynamic table at those coordinates and retain the query
    point, interpolation bracket, axis margins, and source-quality label.
  \item Convert coefficients into dimensional wind or body forces, transform
    into body axes, and add gravity in the same body frame.
  \item Compute propulsion force and mass flow, then transfer propulsion and
    aerodynamic moments to the CG using the declared application points.
  \item Apply actuator maps, limits, and allocation before evaluating the
    derivative. The trim must see the same achieved control that propagation
    will see.
  \item Return named residual components in SI units, for example
    \pathcode{u-dot}, \pathcode{v-dot}, \pathcode{w-dot},
    \pathcode{p-dot}, \pathcode{q-dot}, and \pathcode{r-dot}.
\end{enumerate}

The first estimate from the dimensionalization example is useful here. With
$m=1200$ kg, $V=60$ m/s, $\rho=1.225$ kg/m$^3$, and $S=16$ m$^2$,
$q_\infty=2205$ Pa and $C_L\approx0.334$ supports level lift. The actual
residual still includes body-axis transformation, drag, thrust direction,
gravity projection, moment reference transfer, and the selected controls.
The estimate tells the author whether a failed trim is a numerical problem or
an order-of-magnitude data problem.

\begin{lstlisting}
def plant_residual(state, controls):
    params = vehicle_parameters_at(state["mass"])
    air = relative_air_data(state, params)
    coeffs = aero_tables(air.alpha, air.beta, air.mach, controls)
    forces = dimensional_aero_forces(coeffs, air, params)
    moments = dimensional_aero_moments(coeffs, air, params)
    load = assemble_body_loads(state, controls, forces, moments, params)
    rates = rigid_body_derivative(state, load, params)
    return {name: rates[name] for name in
            ("u-dot", "v-dot", "w-dot", "p-dot", "q-dot", "r-dot")}
\end{lstlisting}

This is intentionally an adapter outline, not a second physics implementation.
The functions must call the same source tables and mass-property provider as
the propagation RHS. If trim and propagation have different force assembly,
the resulting $A/B$ matrix is a derivative of the wrong plant.

\subsection{\texorpdfstring{Generating the $A$ and $B$ matrices}{Generating the A and B matrices}}

Let $f(x,u)$ be the actual plant state-rate evaluator at a source-backed trim
$(x_0,u_0)$. The local continuous model is
\[
  \delta\dot{x} = A\,\delta x + B\,\delta u,
  \qquad
  A = \left.\frac{\partial f}{\partial x}\right|_{x_0,u_0},
  \qquad
  B = \left.\frac{\partial f}{\partial u}\right|_{x_0,u_0}.
\]

For a central finite-difference estimate, each perturbation must be applied
through the same adapter and must remain inside the declared table and actuator
envelopes:
\begin{align}
  A_{ij} &\approx
  \frac{f_i(x_0+\Delta x_j,u_0)-f_i(x_0-\Delta x_j,u_0)}{2\Delta x_j},\\
  B_{ik} &\approx
  \frac{f_i(x_0,u_0+\Delta u_k)-f_i(x_0,u_0-\Delta u_k)}{2\Delta u_k}.
\end{align}

The step is a physical perturbation with a declared unit, not a generic
dimensionless epsilon. Perturbations that cross a table boundary, actuator
limit, branch, or event are rejected or reported separately. A force/moment
Jacobian is not automatically an $A,B$ dynamics pair; it must first be mapped
through mass, inertia, kinematics, and the selected state ordering.

The bridge model may use the six-state local attitude system
\[
  x_{\mathrm{bridge}} =
  [\,\delta\phi,\delta\theta,\delta\psi,p,q,r\,]^{T},
\]
but it remains a kinematic or prescribed-attitude model unless the moments and
inertia are released into the full rigid-body RHS.

\subsection{A slow linearization walkthrough}

The following procedure makes the vehicle parameters visible in every matrix
column. It is deliberately slower than calling a generic numerical Jacobian:

\begin{enumerate}
  \item \textbf{Freeze the accepted trim.} Save $x_0$, $u_0$, mass, CG,
    inertia, atmosphere, wind, table hashes, actuator state, and the exact
    source-quality and uncertainty metadata. The derivative must be local to
    this packet.
  \item \textbf{Declare the order.} For a body-rate design, write the exact
    state and control vectors, for example
    \[
      x=[u,v,w,p,q,r]^T,\qquad
      u_c=[\delta_e,\delta_a,\delta_r,\delta_T]^T.
    \]
    For an attitude bridge, use the six names
    $[\delta\phi,\delta\theta,\delta\psi,p,q,r]$ and do not silently mix
    them with aerodynamic $[\alpha,\beta]$ states.
  \item \textbf{Rebuild dependent quantities.} A perturbed state changes
    relative air velocity, angle coordinates, dynamic pressure, coefficients,
    forces, moments, mass-dependent acceleration, and possibly inertia. The
    evaluator must recompute all of them rather than perturbing only a printed
    telemetry channel.
  \item \textbf{Choose physical steps.} Use a central perturbation that is
    large enough to rise above numerical noise and small enough to remain in
    one smooth table and actuator region. Typical starting points might be
    $0.001$ rad for an angle, $0.01$ rad/s for a body rate, and $0.5$ degrees
    for a control surface, but the source table resolution determines the final
    choice. Record the step and units for every state and control.
  \item \textbf{Evaluate both sides.} For state column $j$, evaluate
    $f(x_0+\Delta x_j,u_0)$ and $f(x_0-\Delta x_j,u_0)$. For control column
    $k$, evaluate $f(x_0,u_0+\Delta u_k)$ and
    $f(x_0,u_0-\Delta u_k)$. Reject a pair that crosses a table boundary,
    actuator clamp, event, or branch.
  \item \textbf{Form the columns.} Divide the difference of each derivative
    component by the physical perturbation. Preserve the state/control names,
    trim values, steps, and source metadata alongside the numeric matrices.
  \item \textbf{Check parameter signatures.} The direct force and moment
    sensitivities should show the expected mass-property scale:
    \[
      \frac{\partial\dot u}{\partial F_x}\approx\frac{1}{m},\qquad
      \frac{\partial\dot p}{\partial M_x}\approx\frac{1}{I_x}.
    \]
    At $m=1200$ kg and $I_x=4200$ kg m$^2$, these are approximately
    $8.33\times10^{-4}$ and $2.38\times10^{-4}$ in SI units. Cross terms from
    $\boldsymbol\omega\times\boldsymbol v$ and
    $\boldsymbol\omega\times(I\boldsymbol\omega)$ are expected away from a
    zero-rate trim.
  \item \textbf{Repeat the step test.} Recompute with half and double steps
    where the table remains smooth. Stable columns should converge within the
    declared numerical tolerance; a large change identifies a table kink,
    branch, saturation, or an unsuitable step.
  \item \textbf{Only then synthesize.} Check dimensions, finite values,
    controllability, conditioning, closed-loop poles, table margin, and
    uncertainty robustness before the matrix enters a controller report.
\end{enumerate}

The implementation artifact is produced by:

\begin{lstlisting}
from taoryx.trim import finite_difference_dynamics_linearization

linearization = finite_difference_dynamics_linearization(
    trim_spec,
    plant_state_derivatives,  # returns one derivative per state name
    trim_result,
    state_step=1.0e-6,
    control_step=1.0e-6,
    metadata={
        "mass-kg": trim_result.state["mass"],
        "inertia-source": "vehicle-provider",
        "aero-source-quality": "estimated",
    },
)
\end{lstlisting}

The helper's steps are fractional scales applied to each variable. For
variables with very different physical units, a vehicle adapter should first
verify the resulting absolute perturbations and, if necessary, expose a
scaled local state or a more specialized perturbation wrapper. The artifact's
state names, control names, trim state, trim controls, $A$, $B$, and metadata
must travel together into profile synthesis.

Mass and aerodynamic estimates remain in context at this stage. If mass or CG
changes materially, repeat the trim and linearization or supply a schedule of
local designs. If the aerodynamic tables are estimated, apply the declared
uncertainty screen and retain table-margin reserve; do not turn a numerically
stable local gain into a claim that the aero model is ground truth.

\subsection{Changing mass and inertia}

The rigid-body model has a generic mass-property seam: an optional
\pathcode{inertia\_provider(state)} is queried by both the derivative and
telemetry paths. With no provider, the registered principal inertias remain
constant. With a provider, the vehicle adapter owns the relationship between
mass, configuration, tank state, payload, center of gravity, and inertia. The
integrator does not infer that inertia is proportional to total mass.

The current declarative extension supplies an explicit dry-mass interpolation
when a source-backed schedule is not yet available:

\begin{lstlisting}
*runtime status vehicle inertia-mass-model=linear-dry-mass
  dry-mass-kg=100
  inertia-x=12 inertia-y=18 inertia-z=25
  inertia-x-dry=9 inertia-y-dry=14 inertia-z-dry=20
*runtime lqr attitude update=mass ...
\end{lstlisting}

The reference inertias are the values at the initial mass. The runtime
interpolates to the three explicitly supplied dry-mass inertias, clamps beyond
the declared interval, and publishes
\pathcode{mass\_kg}, \pathcode{propellant\_mass\_kg}, and
\pathcode{inertia\_{x,y,z}\_kg\_m2}. This is a declared approximation, not a
claim about tank geometry or a hidden mass-scaling law. A vehicle with center
of gravity migration, jettison, or an off-diagonal inertia tensor needs a
vehicle-specific provider instead.

For the direct-moment attitude bridge, the scheduled input block is
\[
  B_{\omega M}(m) =
  \operatorname{diag}\!\left(\frac{1}{I_x(m)},
  \frac{1}{I_y(m)},\frac{1}{I_z(m)}\right).
\]
\pathcode{update=mass} rebuilds the cached gain when the supplied mass or
inertia leaves its tolerance band. This accounts for changing rotational
authority; it does not pretend that an attitude-only six-state bridge contains
the translational effects of mass. To include those effects, generate a full
source-trim $A/B$ pair whose state-rate evaluator includes the current mass,
inertia, propulsion, and aerodynamic model. A fixed \pathcode{b-table} cannot
claim mass scheduling.

\subsection{Scaled continuous-time LQR}

The continuous-time LQR minimizes
\[
  J = \int_0^\infty
  \left(x^{T}Qx+u^{T}Ru\right)\,dt
\]
for
\[
  \dot x=Ax+Bu, \qquad u=-Kx.
\]
The gain is obtained from the positive semidefinite solution $P$ of the
continuous algebraic Riccati equation:
\[
  A^{T}P+PA-PBR^{-1}B^{T}P+Q=0,
  \qquad K=R^{-1}B^{T}P.
\]

Raw SI values can make $Q$ and $R$ ill-conditioned across vehicles. Define
diagonal state and control scales $S_x$ and $S_u$ by
\[
  \widetilde{x}=S_x x, \qquad \widetilde{u}=S_u u.
\]
The normalized system is
\[
  \widetilde A=S_x A S_x^{-1}, \qquad
  \widetilde B=S_x B S_u^{-1},
\]
and the physical gain is recovered with
\[
  K_{\mathrm{physical}}=S_u^{-1}\widetilde K S_x.
\]

The current scale contract derives angle, rate, and moment scales from the
registered vehicle parameters, actuator authority, and table extents. These
scales improve conditioning; they are not permission to command beyond a table
or actuator limit.

\subsection{Automatic profiles and fail-closed selection}

Every vehicle has three named profiles:

\begin{description}[leftmargin=1.1in,labelwidth=0.9in,style=sameline]
  \item[gentle] Higher control penalty and slower response for early plant
    development or a narrow local envelope.
  \item[standard] The default development profile, balanced for local recovery
    and nominal actuator use.
  \item[aggressive] Faster response for cases with demonstrated authority and
    sufficient table margin; it is never a default claim.
\end{description}

\pathcode{python tools/auto_tune_controllers.py --vehicle all} synthesizes the
profiles from the registry. The default report is explicitly
\texttt{runtime-inertia-attitude-bridge, synthesis-only} when no vehicle
source-trim Jacobian and maneuver evaluator are supplied. That report proves
matrix conditioning and closed-loop pole placement only.

A source-backed candidate must also pass the supplied limits for:

\begin{itemize}
  \item maximum closed-loop pole real part;
  \item minimum normalized table margin;
  \item control saturation fraction;
  \item absolute control-rate maximum;
  \item tracking error, when a maneuver evaluator is available.
\end{itemize}

The tuner serializes the scales, their provenance, table axis limits, weights,
closed-loop eigenvalues, condition number, metrics, violations, score, and
status. A stable pole set with table extrapolation or persistent saturation is
not a safe candidate. The winning candidate is copied into a fixed problem
configuration and rerun through closure, convergence, source parity, and
mission gates; auto-tuning alone is not mission evidence.

\subsection{From a bounded trim to a source-informed controller}

The minimum reusable adapter has two distinct evaluators. The trim evaluator
returns named residuals such as force and moment balance. The dynamics
evaluator returns named state derivatives in exactly the same state ordering;
it must not return residual forces or moments in their place. A typical source
informed tuning path is:

\begin{lstlisting}
from taoryx.controller_autotune import AutoTuneLimits, auto_tune_lqr_profiles
from taoryx.runtime import LqrUncertaintySpec
from taoryx.trim import finite_difference_dynamics_linearization

trim = solve_trim(trim_spec, plant_residual)
linearization = finite_difference_dynamics_linearization(
    trim_spec,
    plant_state_derivatives,
    trim,
    state_step=1.0e-6,
    control_step=1.0e-6,
    metadata={"source": "vehicle-adapter", "mass-kg": trim.state["mass"]},
)
report = auto_tune_lqr_profiles(
    "new-vehicle",
    linearization=linearization,
    limits=AutoTuneLimits(
        derivative_uncertainty=LqrUncertaintySpec(
            a_fraction=0.15, b_fraction=0.20, samples=25,
        ),
    ),
    evaluator=run_bounded_maneuver,
)
\end{lstlisting}

The resulting report is labelled
\texttt{source-trim-dynamics-linearization}, retains the trim state and
control names, records the perturbation steps and metadata, and can gate on
closed-loop poles, table margin, saturation, rate, tracking error, and sampled
uncertainty. The source-trim Jacobian is local: repeat the trim and design at
each materially different mass, center of gravity, or flight condition, or
provide an explicit schedule of local designs.

\subsection{Estimated aerodynamic data and design reserve}

Aerodynamic tables and derivatives are estimates unless their provenance says
otherwise. The controller workflow therefore separates three claims:

\begin{enumerate}
  \item table validity: the query remains inside the hard source-table
    boundary;
  \item operational reserve: the normalized table margin remains above the
    declared reserve after subtracting the model uncertainty fraction; and
  \item local robustness: the fixed gain remains stable for the deterministic
    bounded perturbations used by the uncertainty screen.
\end{enumerate}

For example:

\begin{lstlisting}
*runtime status vehicle aero-uncertainty-fraction=0.15
  aero-source-quality=estimated
*runtime lqr attitude a-uncertainty-fraction=0.15
  b-uncertainty-fraction=0.20 uncertainty-samples=25
  uncertainty-policy=fail-closed ...
\end{lstlisting}

The sampled screen is evidence about the declared local envelope, not a formal
robust-control proof, confidence interval, or correction to the aerodynamic
table. Uncertainty values must be traceable to source scatter, digitization,
validation data, or an explicitly recorded engineering assumption. A stable
gain with an estimated aerodynamic model is still an estimated-model result.

\section{Trajectory objectives and scoring}
\label{sec:scoring}

Trajectory evaluation consumes canonical SI telemetry and declared events. It
does not infer units from a plot label and it does not score a reference command
as if it were the actual state. Every objective declares an identifier, kind,
channel, target or bound, tolerance, unit, comparison, phase, weight, and
severity.

\subsection{Objective types}

The common evaluator supports:

\begin{itemize}
  \item state targets: position, altitude, speed, heading, alpha, beta, or
    body rates;
  \item bounded channels: Mach, dynamic pressure, heating proxy, load factor,
    table margin, actuator position, or actuator rate;
  \item events: waypoint capture, segment transition, separation, mode handoff,
    ground contact, or declared termination;
  \item trajectory properties: cross-track error, route completion, return to
    start, and recovery after a disturbance;
  \item invariants: continuity, mass balance, quaternion norm, force/moment
    closure, and step-size convergence.
\end{itemize}

A waypoint is an event plus a measured objective. Its result must retain the
actual position, target position, arrival time, distance error, tolerance, and
slack. For a rectangle or figure-eight, score each corner and the whole route;
do not replace the route with only the final point.

\subsection{Per-objective result}

For an absolute-error objective, the evaluator computes
\[
  e=|y-y^{\star}|, \qquad
  s=\tau-e, \qquad
  e_{\mathrm{norm}}=\frac{e}{\tau}, \qquad
  q=\max(0,1-e_{\mathrm{norm}}),
\]
where $\tau$ is the declared tolerance. The raw result retains the unit and
time or event location:

\begin{lstlisting}
{
  "id": "waypoint-2",
  "status": "pass",
  "actual": 301.4,
  "target": 300.0,
  "error": 1.4,
  "tolerance": 2.0,
  "slack": 0.6,
  "normalized_error": 0.7,
  "unit": "m",
  "time_s": 42.0,
  "weight": 1.0,
  "severity": "required"
}
\end{lstlisting}

For a band, the error is the distance outside the band. For an event, the
result records whether the declared event occurred and its event audit. Missing
or unitless telemetry is \texttt{blocked}, not zero.

\subsection{Trajectory rollup}

The rollup contains two layers.

\paragraph{Validity gates.} Required objectives, source parity, continuity,
convergence, equation closure, full table-envelope validity, actuator
feasibility, and declared mission completion are gates. Any required failure or
missing channel blocks the claim. Advisory objectives can produce a
\texttt{diagnostic} status but cannot hide a required failure.

\paragraph{Quality score.} After the declared gate policy is applied, an
informative composite may be calculated as
\[
  \operatorname{score}=100\frac{\sum_i w_i q_i}{\sum_i w_i},
  \qquad 0\leq q_i\leq1.
\]
The rollup must also expose the worst required normalized error and every
negative slack. A trajectory that barely stays inside a 300 m hard limit must
not receive the same quality score as one that closes to 2 m.

The resource and validity metrics are part of the rollup, with units and
sampling definitions:

{\small
\begin{longtable}{p{0.35\linewidth}p{0.2\linewidth}p{0.35\linewidth}}
\toprule
\textbf{Metric} & \textbf{Unit} & \textbf{Meaning}\\
\midrule
\endhead
\texttt{table\_margin\_min\_normalized} & dimensionless &
  Minimum distance to any queried table edge divided by that axis span; zero is
  the hard validity boundary.\\
\texttt{table\_margin\_average\_normalized} & dimensionless &
  Time-sample average of the normalized margin.\\
\texttt{control\_saturation\_fraction} & dimensionless &
  Fraction of samples with any saturation flag active.\\
\texttt{control\_saturation\_average} / \texttt{max\_abs} & dimensionless &
  Aggregate saturation severity, with the per-channel breakdown retained.\\
\texttt{control\_derivative\_abs\_average} & command unit/s &
  Mean absolute command derivative over saved samples.\\
\texttt{control\_derivative\_abs\_max} & command unit/s &
  Maximum absolute command derivative; report each control separately.\\
\bottomrule
\end{longtable}
}

The final composite dictionary contains status, score, gate score, quality
score when available, required counts, worst normalized error, all objective
records, termination information, and the scenario-contract hash. The source
file is \pathcode{src/taoryx/objectives.py}; the evidence packet also carries
\pathcode{score_definition.json} and the metric dictionary so that a reviewer
can recompute the rollup.

\subsection{Family-specific objective phases}

Objectives must follow the mechanics of the vehicle:

\begin{itemize}
  \item \textbf{Powered fixed-wing:} trim hold, altitude capture, heading or
    bank corner, speed recovery, waypoint sequence, and return corridor.
  \item \textbf{Unpowered glider or X-15 segment:} release state, energy or
    descent corridor, alpha/beta and table margins, bank reversal, and valid
    ground-contact or termination event.
  \item \textbf{Quadrotor:} takeoff, hover, north/east waypoint, yaw step,
    disturbance recovery, return, and landing. Include rotor commands and
    achieved speeds in the objective evidence.
\end{itemize}

Event markers belong to the owning family and phase. The showcase plot should
show the target waypoint or objective, capture time, active phase, and the
change in controller mode. It must not place unrelated family events on a
shared timeline.

\section{Verification sequence}
\label{sec:verification}

The common harness follows the same sequence for every 3-DOF, bridge, or
6-DOF vehicle. A later stage cannot promote a blocked earlier stage.

\begin{enumerate}
  \item \textbf{Table ingestion and axis order:} parse every bound table,
    verify cardinality, independent-variable order, coefficient names, and
    source hash.
  \item \textbf{Units and reference geometry:} check native-to-SI conversion,
    area/span/chord, mass, CG, and moment reference.
  \item \textbf{Body/world frames:} verify basis orientation, handedness,
    active/passive transforms, gravity, and wind direction.
  \item \textbf{Quaternion identity:} confirm ordering, identity rotation,
    normalization, determinant, and round trips.
  \item \textbf{Alpha/beta conventions:} derive angles from relative velocity,
    exercise sign probes, zero-speed behavior, and branch limits.
  \item \textbf{Coefficient lookup and margins:} compare source evaluator and
    table evaluator, record exact query coordinates, and fail closed on
    extrapolation.
  \item \textbf{Force/moment dimensionalization:} compare coefficient-based
    and dimensional wrenches with separate axes and units.
  \item \textbf{Mass, CG, and inertia:} verify positivity, symmetry,
    continuity, eigenvalues, and parallel-axis updates.
  \item \textbf{Static wrench or trim closure:} solve or import a bounded trim;
    retain unscaled force and moment residuals.
  \item \textbf{Short bounded propagation:} run a small deterministic case with
    complete telemetry and a declared termination policy.
  \item \textbf{Integration convergence:} compare $\Delta t$, $\Delta t/2$,
    $\Delta t/4$, and an adaptive reference on smooth intervals.
  \item \textbf{Controller and mission tests:} only now score waypoint,
    altitude, recovery, disturbance, and whole-trajectory objectives.
\end{enumerate}

The standard plant report exposes each stage as \texttt{pass}, \texttt{diagnostic},
\texttt{blocked}, or \texttt{not\_run}. For each four-family report the first
eleven stages are plant/convention evidence; controller and mission claims are
not implied by a plant-golden result.

\subsection{The validation ladder for a new vehicle}

The twelve checks above are the core harness, but a new vehicle also needs a
clear promotion path. The rule is monotonic: a later green trajectory cannot
promote an earlier blocked data, frame, table, or trim check. Each stage must
leave an inspectable artifact for the next stage.

\begin{center}
\begin{tabularx}{0.96\linewidth}{>{\bfseries}p{0.19\linewidth}p{0.27\linewidth}Y}
\toprule
Stage & Question & Required evidence \\
\midrule
Registry and data & Can the model be found and loaded? & SI metadata, family contract, table bindings, source provenance/hash, generated problem profile, onboarding report. \\
Grammar and lowering & Does the declaration mean what the author intended? & Grammar diagnostics, segment lint/build, resolved manifest, and transition audit. \\
Convention firewall & Are frames, table axes, signs, and controls correct? & Table inspection, nominal/boundary queries, frame and quaternion tests, and signed control probes. \\
Plant validity & Does the model satisfy its equations? & Initial-condition audit, trim residuals, dimensional wrench checks, independent closure, and bounded propagation. \\
Numerical quality & Is the result reproducible and time-step credible? & $\Delta t$, $\Delta t/2$, $\Delta t/4$, adaptive comparison, event-aware exclusions, and output hashes. \\
Controller validity & Does the controller stabilize the local plant? & Named $A/B$ provenance, controllability, LQR poles, uncertainty screen, and actuator/slew/allocation telemetry. \\
Segment validity & Do entry, handoff, and exit contracts hold? & Inherited/reset state audit, controller reset events, goals, and termination evidence. \\
Checkpoint mission & Does the complete scenario meet bounded objectives? & Generated fidelity packet, objective gates, quality metrics, plots, and claim status. \\
\bottomrule
\end{tabularx}
\end{center}

This separates four questions that are often accidentally combined: \emph{is
the model registered}, \emph{is the plant physically and numerically valid},
\emph{is the controller valid for that plant}, and \emph{does this particular
scenario meet its bounded objectives}. A model can be ready for the first
question and still be blocked at the other three.

\subsection{Initial vehicle addition and data firewall}

The initial addition process is metadata-first. Preserve the source bundle and
its provenance, add one SI entry to
\pathcode{verification/vehicle_models.yaml}, select a family contract, bind
the tables, add a generated-problem profile and golden catalog case, then add
the trim, controller, control-direction, and LQR profile records. The focused
validator checks completeness of that path; it does not claim that a trim or
mission has passed.

\begin{lstlisting}
python tools/validate_vehicle_onboarding.py --vehicle new_vehicle --strict
python tools/dev.py generate-problems
python tools/dev.py vehicles
taoryx table inspect path/to/vehicle.tbl --html build/table-explorer.html
python tools/dev.py control-directions
python tools/dev.py trim-vehicles
\end{lstlisting}

The data packet supplied before trim should answer all of the following:

\begin{enumerate}
  \item What is the exact state and control order, and which variables are
    body, wind, inertial, or world quantities?
  \item What are the SI conversions, reference area, span, chord, mass, CG,
    moment reference, and declared operating envelope?
  \item What is the table independent-variable order, cardinality, monotonicity,
    interpolation policy, valid bound, coefficient frame, and source hash?
  \item What happens at a nominal point, every relevant boundary, and a
    deliberately out-of-range query? Extrapolation must be a declared
    diagnostic or block, not a silent continuation.
  \item Are mass, CG, and inertia positive, symmetric, continuous, and
    supplied by the intended provider at every operating point?
  \item Which aerodynamic quantities are measured, executable-source values,
    digitized, estimated, or assumed, and what uncertainty bound applies?
\end{enumerate}

The table firewall must preserve source orientation and signs. Never reverse an
axis, negate a coefficient, or change a moment reference solely because the
first trajectory looks wrong. Compare the source evaluator with the table
evaluator at recorded coordinates and report the table margin, source/table
error, and worst query. A force or moment coefficient is not a body-axis force
or moment until its wind-to-body transform and reference transfer are explicit.

The control convention harness uses the same baseline for the plus and minus
perturbations of one control. It records the signed body-force or body-moment
response, antisymmetry error, expected source sign, and achieved bounded
command. Representative assertions include the B747 elevator's body-$y$
moment sign, the X8 collective and differential elevon separation of body-$y$
and body-$x$ response, the Hummingbird rotor-speed body-$z$ force sign, and the
X-15 stabilator/rudder body-axis signs. These tests are a firewall against a
wrong table orientation or actuator convention; they are not controller tuning
tests.

\subsection{LQR stability and controller gates}

An LQR result is promotable only when the plant and the controller agree on
names, units, and directions. The controller artifact must include the trim
identifier, exact state/control ordering, perturbation sizes, $A$ and $B$
matrices, $Q$ and $R$, scaling, controllability result, nominal closed-loop
eigenvalues, actuator limits, slew limits, allocator status, table margins,
and the uncertainty cases used for the design.

The minimum pole gate is
\[
  \max_i \operatorname{Re}\lambda_i(A-BK)<0,
\]
with a declared numerical margin and no unstable allowed uncertainty corner.
Check the controllability rank before solving, and check the saturated and
rate-limited command path after solving. A stable pole set with reversed
control signs, table extrapolation, or persistent saturation is a failed
controller gate, not a qualified design.

For changing mass properties, the runtime may use a declared generic inertia
provider and a mass-scheduled attitude update. That provider updates the
attitude inertia seen by the bridge; it does not infer inertia from mass, and it
does not replace a fresh plant-bound $A/B$ design when mass changes CG,
translation, propulsion, or aerodynamic derivatives. For estimated
aerodynamics, evaluate the declared parameter corners or perturbation model
and fail closed when the worst-case pole, table margin, or actuator reserve is
not acceptable.

\subsection{Parametric validity trajectories}

Checkpoint trajectories should be generated from machine-readable catalogs,
not chosen after inspecting which run looks attractive. The principal sources
are:

\begin{itemize}
  \item \pathcode{verification/fidelity_ladder.yaml};
  \item \pathcode{verification/long_validation_trajectories.yaml};
  \item \pathcode{verification/controller_missions.yaml};
  \item \pathcode{verification/vehicle_maneuver_matrix.yaml}; and
  \item \pathcode{verification/segmentation_catalog.yaml}.
\end{itemize}

These catalogs parameterize the
vehicle, source problem and tables, duration, maximum steps, integration
factors, envelope bounds, objective targets, and termination policy.

The generated fidelity sequence is:

\begin{enumerate}
  \item source or static trim hold;
  \item historical or source 3-DOF anchor;
  \item derived kinematic 3+3 bridge, with translational parity and an
    attitude sidecar;
  \item short source-anchored rigid-body 6-DOF propagation;
  \item $\Delta t$, $\Delta t/2$, $\Delta t/4$, and adaptive convergence;
  \item source differential and independent closure;
  \item bounded recovery, disturbance, or control maneuver; and
  \item long validation and checkpoint mission.
\end{enumerate}

\begin{lstlisting}
python tools/dev.py fidelity-packet
python tools/dev.py maneuver-matrix
\end{lstlisting}

These commands produce the UUID-scoped ladder packet and bound maneuver
catalog. A valid packet contains the resolved inputs and hashes, initial
condition audit, event timeline, closure metrics, convergence report, source
differential result, objectives, and plots. The required artifact set is
defined by the long-validation catalog; a trajectory that merely completes is
not sufficient evidence.

\subsection{Segment entry, transition, and exit validation}

The segmentation catalog is an orchestration contract around a source problem.
Each ordered segment names its entry condition, exit condition and target or
action, initial or inherited state, controller, actuator binding, table and
environment dependencies, goal, and transition policy. The compiler rejects
duplicate identifiers, missing targets, accidental cycles, invalid
\texttt{stop}/\texttt{goto} forms, missing source segments, and source
segments without an integration block.

Validate a segment in three separate moments:

\begin{enumerate}
  \item \textbf{Entry:} confirm the initial or inherited state is complete and
    unit-bearing; position, velocity, attitude, rates, mass, CG, and
    quaternion are finite and valid; tables are in range; and the active
    controller and actuator binding agree with the segment contract.
  \item \textbf{Transition:} verify continuity of every declared field, or
    verify an explicit reset, impulse, separation, or mass change. Record the
    event in the transition audit and reset the controller at the first
    segment and at every handoff.
  \item \textbf{Exit:} prove that the declared condition was observed, the
    target or action exists, the final state remains valid, and the termination
    reason is intentional rather than a timeout, table failure, safety clamp,
    or numerical exception.
\end{enumerate}

Run \texttt{segment-lint}, \texttt{segment-build}, and \texttt{segment-run},
then inspect the generated manifest and transition audit. Compilation proves
composition and syntax only; it does not prove plant closure, convergence,
envelope validity, or controller quality.

\subsection{Checkpoint quality and promotion}

Keep hard validity gates separate from quality metrics. Required gates include
completion under the declared termination policy, initial-condition validity,
table margin, source differential where applicable, force/moment closure,
time-step convergence, mass and quaternion continuity, envelope limits,
saturation and rate limits, segment events, and controller poles. Quality
metrics may then report waypoint capture error, drift, recovery time, margin
reserve, control effort, and smoothness.

The status in \pathcode{verification/staged_completion_matrix.yaml} is bounded
by the lowest failing tier. Use \texttt{pass}, \texttt{pass\_local}, or the
catalog's more limited directional/source-backed status only when the named
evidence exists; use \texttt{blocked} when a required gate is absent or fails.
The final report must name the vehicle, fidelity tier, source status, model
uncertainty, controller profile, scenario, and limitations. No plot, score, or
stable eigenvalue may silently substitute for missing source data or a failed
validity gate.

\subsection{Independent equation closure}

Closure must be computed from the saved right-hand-side derivative and an
independent force/moment recomputation, with every vector carrying an explicit
frame. For body-frame translation and rotation:
\begin{align}
  r_F &= m\left(\dot{\boldsymbol v}_B+
    \boldsymbol\omega_B\times\boldsymbol v_B\right)
    -\left(\boldsymbol F_{A,B}+\boldsymbol F_{P,B}+\boldsymbol F_{G,B}\right),\\
  r_M &= I\dot{\boldsymbol\omega}_B+
    \boldsymbol\omega_B\times(I\boldsymbol\omega_B)-\boldsymbol M_B.
\end{align}
Normalized values use declared scales, for example
\[
  \epsilon_F=\frac{\lVert r_F\rVert}{\max(mg,\lVert F_B\rVert)},
  \qquad
  \epsilon_M=\frac{\lVert r_M\rVert}{\max(1\ {\rm N\,m},\lVert M_B\rVert)}.
\]

If a separate finite-difference closure is also emitted, it must be labelled
separately from the algebraic RHS closure. Events such as ignition, separation,
control clamps, and ground contact are one-sided samples or explicitly excluded
with an event reason; they must not be silently folded into a p99 statistic.

\subsection{Source differential evidence}

Before a long trajectory, sample the declared source envelope and compare the
source evaluator with the \taoryx{} table or wrench evaluator. Publish sample
count, maximum absolute and relative error, RMS error, worst input state, and
source/table hashes. Exact executable transcodes can target near floating-point
agreement; digitized or rounded public data require a documented uncertainty
bound. A trajectory that looks plausible does not replace this test.

\subsection{Convergence and reproducibility}

Each golden case records the integrator, step sizes, output rate, tolerances,
software commit, Python/dependency environment, problem and table hashes,
controller configuration, termination policy, and all output hashes. The
evidence packet is UUID-scoped and contains raw telemetry, events, objectives,
closure, convergence, source parity, plots, and the manifest. Generated
showcases are presentation products; they do not replace the machine-readable
packet.

\section{A repeatable authoring workflow}

The shortest safe route for a new vehicle is:

\begin{lstlisting}
# 1. Diagnose the metadata and source contract.
python tools/validate_vehicle_onboarding.py --vehicle new_vehicle

# 2. Generate ordinary problem files from catalog metadata.
python tools/dev.py generate-problems

# 3. Verify tables, provenance, signed controls, and generated products.
python tools/dev.py vehicles
python tools/dev.py control-directions

# 4. Solve the bounded source-backed trim and synthesize profiles.
python tools/dev.py trim-vehicles
python tools/auto_tune_controllers.py --vehicle new_vehicle

# 5. Run only the selected family view while developing.
python tools/dev.py test-x8

# 6. Build the composite extension PDF and complete repository checks.
python tools/dev.py taoryx-extension-pdf
python tools/dev.py check
\end{lstlisting}

The exact family command changes with the vehicle; the workflow does not. A new
family is justified only when its frame, dynamics, actuator, or control
semantics cannot be represented by an existing family contract. New problem
files should be generated from metadata rather than copied with constants
embedded in each test.

\section{Claim vocabulary and release checklist}

The evidence packet uses explicit claim labels:

\begin{center}
\begin{tabularx}{0.94\linewidth}{>{\bfseries}p{0.22\linewidth}Y}
\toprule
manual-bounded & Supported by the reconstructed 1995 manual within its documented scope.\\
taoryx-extension & Successor behavior deliberately outside the historical claim.\\
scaffolded & The source or contract exists, but the complete runtime or independent oracle is unfinished.\\
verified & The declared source, runtime behavior, and independent assertions pass the stated gates.\\
\bottomrule
\end{tabularx}
\end{center}

Before a controller-mission claim is released, confirm:

\begin{itemize}
  \item the selected vehicle is onboarding \texttt{ready};
  \item tables, units, frames, quaternion, angle signs, and control directions
    pass the convention firewall;
  \item trim and source differential artifacts are packaged;
  \item $A$ and $B$ provenance, perturbation sizes, eigenvalues, and scaling
    are recorded;
  \item all required objectives have unit-bearing results and nonnegative
    slack, or the run is explicitly failed/blocked;
  \item table margins, saturation, command derivatives, closure, continuity,
    convergence, and termination are reported;
  \item the final claim names the fidelity tier, source status, family, and
    limitations; and
  \item no score, plot, or green test silently substitutes for missing source
    data or a failed validity gate.
\end{itemize}

\warningbox{\textbf{Do not overclaim.} A stable LQR pole set, a finite trajectory,
or a high quality score does not establish global aerodynamic validity,
historical TAOS compatibility, or engineering certification. Those claims
require their own source authority, independent evidence, and release gate.}

\section{Traceability map}

This guide is an entry point, not a second source of truth. The executable
contracts and tests remain authoritative for current successor behavior:

\begin{center}
\begin{tabularx}{0.94\linewidth}{>{\raggedright\arraybackslash}p{0.34\linewidth}Y}
\toprule
\textbf{Topic} & \textbf{Repository authority}\\
\midrule
Fidelity tiers and parity & \pathcode{src/taoryx/scenario_contract.py}, \pathcode{verification/fidelity_ladder.yaml}\\
Vehicle onboarding & \pathcode{src/taoryx/vehicle_onboarding.py}, \pathcode{docs/plan/vehicle-onboarding.md}\\
Trim and local derivatives & \pathcode{src/taoryx/trim.py}, \pathcode{verification/trim_specs.yaml}\\
LQR solve and runtime contract & \pathcode{src/taoryx/runtime/lqr.py}, \pathcode{docs/extensions/lqr.md}\\
Profile auto-tuning & \pathcode{src/taoryx/controller_autotune.py}, \pathcode{tools/auto_tune_controllers.py}\\
Objective scoring & \pathcode{src/taoryx/objectives.py}, \pathcode{docs/plan/scored-fidelity-and-mission-evidence-v1.md}\\
Plant firewall & \pathcode{tests/e2e/support/golden_plants.py}, \pathcode{tests/e2e/test_golden_vehicle_plants.py}\\
Build entry point & \pathcode{tools/dev.py} \pathcode{taoryx-extension-pdf}, \pathcode{docs/BUILDING.md}\\
\bottomrule
\end{tabularx}
\end{center}

\end{document}
